\scanpage{94}
de $\Pi_0$, $\Pi'_0=\operatorname{Im}(\Pi_1(\widetilde{\mathcal C},x'_0)\mathbin{-}\Pi_1(\mathcal C,x_0))$, de sorte que \deleted{$\unsure{\widetilde{\mathcal C}|\widetilde E}$}
\[
(47)\qquad E_0\simeq\Pi_0/\Pi'_0;
\]
la condition d'existence de $\widetilde{\Theta}_0$ revient à dire
\[
\Theta(\Pi'_0)\text{ conjugué à }\Pi'_0\text{ dans }\Pi_0
\]
\margin{en effet, il faudra que $\Theta(\Pi'_0)$ soit $\operatorname{Im}(\Pi_1(\widetilde{\mathcal C},x'_1)\unsure{\mathbin{-}|\to}\underbrace{\Pi_1(\mathcal C,x_0)}_{\Pi_0})$, où $x'_1=\widetilde{\Theta}(x'_0)$, Or si $x'_1=g(x'_0)$, où $g\in\Pi_0$, l'Im en question vaut $\operatorname{int}(g)(\Pi'_0)$}
et les choix possibles pour remonter $\Theta$ en $\widetilde{\Theta}_0$, donc pour remonter $\Theta$ en $\widetilde{\Theta}$, correspondent aux éléments \inserted{$\unsure{x'_1|x''}$} de $E_0$ \deleted{tels que} \inserted{de la} forme $g(x'_0)$ avec $g\in\Pi_0$ satisfaisant
\[
(48)\quad\unsure{;|:}\qquad \operatorname{int}g(\Pi'_0)=\Theta(\Pi'_0)
\]
i.e. tels que
\[
g\in\operatorname{Transp}_{\Pi_0}(\Pi'_0,\Theta(\Pi'_0))
\]
\margin{Cette condition ne dépend que de l'endom. \uline{ext.} défini par $\Theta$ sur $\Pi'_0$ --- mais la façon de remonter $\Theta$ en $\widetilde{\Theta}$, $=$ comme endom. ext., en dépend, \unsure{aux de l'} \uline{pas} \unsure{un|en} endom. ord. \unsure{défini} par $\Theta$}
\margin{\deleted{\unsure{pts} \illegible{2}}\quad\deleted{\unsure{Si} $\unsure{\widetilde{\mathcal C}}$ \unsure{correspond} \illegible{3} \unsure{on a} $\unsure{\Theta'_0}$ \illegible{3} $\unsure{\Theta}$ \illegible{2} $\unsure{\widetilde{\Theta}}$ \unsure{tel que} $\unsure{E_0}\ \unsure{\exists\gamma\in\Pi_0}$ \unsure{avec} $(48\,\mathrm{bis}):\ \Theta(\Pi'_0)=\gamma(\Pi'_0)$ dans $\Pi'_0$.}}
donc les $\Theta$ possibles correspondent aux éléments de
\[
(49)\qquad \operatorname{Transp}_{\Pi_0}(\Pi'_0,\Theta(\Pi'_0))/\Pi'_0
\]
\margin{espace homogène à droite ($\Pi'_0$ opère à droite sur le transporteur, via $g\longmapsto g\alpha$ pour $\alpha\in\Pi_0$}
\margin{Bien sûr, les façons de remonter $\Theta$: $\Theta'$ de $\mathcal C'\subset\widetilde{\mathcal C}$ correspondent aux $\widetilde{\Theta}$ qui rendent bien ainsi $\mathcal C'\subset\widetilde{\mathcal C}$ invariant.}

Ainsi, on voit que dans le cas \deleted{$\unsure{\mathcal C|\widetilde{\mathcal C}}$}
\[
\mathcal C=\widehat{\Pi}_1(X\deleted{_{\unsure{\overline{k}|\overline K}}},I),\qquad I\subset X_{\unsure{k|K}}
\]
quand $\mathcal C'$ \deleted{en premier} est de la forme $\widehat{\Pi}_1(X'\deleted{_{\unsure{\overline{k}|\overline K}}},I')$, où $X'$ est un rev. étale (fini) de $X$, \inserted{avec $I'\unsure{\mathbin{-}|\to}I$}, que la condition qu'un \inserted{\unsure{endom}} $\Theta$ \inserted{de $\mathcal C$} se remonte en un autom de $\mathcal C'$ est une condition algébrique des plus
\margin{93}
